\documentclass[../main.tex]{subfiles}

\begin{document}

\section{Khảo sát hiện trạng}
Sinh viên mô tả hiện trạng, quy trình nghiệp vụ hoặc bài toán thực tế liên quan đến đề tài.

\section{Phân tích yêu cầu}
Sinh viên trình bày yêu cầu chức năng, yêu cầu phi chức năng và các ràng buộc của hệ thống hoặc giải pháp.

\section{Phương pháp thực hiện}
Sinh viên mô tả phương pháp, quy trình, công cụ và dữ liệu sử dụng trong quá trình thực hiện đề tài.

\section{Kế hoạch thực hiện}
Sinh viên trình bày các giai đoạn công việc, mốc thời gian và sản phẩm dự kiến của từng giai đoạn.
Ví dụ, kế hoạch thực hiện có thể được trình bày như Bảng~\ref{tab:ke-hoach}.

\begin{table}[H]
\centering
\caption{Ví dụ bảng kế hoạch thực hiện}
\label{tab:ke-hoach}
\begin{tabular}{|c|p{5cm}|p{5cm}|}
\hline
\textbf{STT} & \textbf{Công việc} & \textbf{Kết quả dự kiến} \\
\hline
1 & Khảo sát và thu thập yêu cầu & Danh sách yêu cầu \\
\hline
2 & Phân tích và thiết kế & Tài liệu thiết kế \\
\hline
3 & Triển khai và kiểm thử & Sản phẩm và kết quả kiểm thử \\
\hline
\end{tabular}
\end{table}

\end{document}
